% AIC 201 - Assignment 3 (CCP): AIDRA technical report
% IEEE conference format. Compile: pdflatex (twice) or latexmk -pdf
% Expanded body text + tables; with one architecture PDF expect ~5--6 pages.
\documentclass[conference]{IEEEtran}
\IEEEoverridecommandlockouts
\usepackage[T1]{fontenc}
\usepackage{times}
\usepackage{graphicx}
\usepackage{amsmath,amssymb}
\usepackage{booktabs}
\usepackage{caption}
\usepackage{balance}
\usepackage{url}
\usepackage[hidelinks]{hyperref}
\usepackage{microtype}
\usepackage{algorithm}
\usepackage{algpseudocode}

\setlength{\textfloatsep}{8pt plus 2pt minus 2pt}
\setlength{\floatsep}{8pt plus 2pt minus 2pt}
\setlength{\intextsep}{8pt plus 2pt minus 2pt}

\newcommand{\Aidra}{AIDRA}

\begin{document}

\title{\Aidra: Integrated Search, CSP, Machine Learning, and Fuzzy Routing\\
for Disaster Response Simulation}

\author{\IEEEauthorblockN{Student Name~1,\, Student Name~2}
\IEEEauthorblockA{Department of Computer Science\\
Bahria University Islamabad Campus (example)\\
Email: replace.with.your@students.bahria.edu.pk}}

\maketitle

\begin{abstract}
This report documents the Adaptive Intelligent Disaster Response Agent (\Aidra), a
full-stack intelligent system built for the Complex Computing Problem (CCP) in
Artificial Intelligence (AIC~201). The implementation targets Course Learning
Outcome CLO3: appraising real-world problems for machine learning based solutions,
here in the context of \emph{integrated} AI rather than an isolated classifier.
\Aidra{} couples classical graph search (breadth-first, depth-first, greedy
best-first, and A*), constraint satisfaction with backtracking and heuristics,
optional local search post-processing, a supervised three-class risk estimation
pipeline (k-nearest neighbors, Gaussian Naive Bayes, and a one-hidden-layer
multilayer perceptron trained on synthetic data), and Mamdani-style fuzzy rules that
reshape planner costs from hazard, urgency, and uncertainty signals. The deployed
application simulates an $18\times 18$ dynamic grid, five victims, two ambulances
(capacity two each), a rescue team, medical centers, survival decay, shock and
hazard events, and explicit replanning with decision logging. We give a formal
problem summary, architecture, algorithmic detail, tabular summaries of features
and parameters, evaluation metrics aligned with the live UI, interpretation
guidance, and submission artifacts (GitHub, two-minute LinkedIn demo).
\end{abstract}

\begin{IEEEkeywords}
disaster response, A*, constraint satisfaction, machine learning, fuzzy logic, replanning, simulation, CLO3
\end{IEEEkeywords}

\section{Introduction}

\subsection{Motivation}
Urban disaster response requires coordinating scarce mobile resources under partial
information, heterogeneous victim severity, and a terrain that may degrade during
the incident---roads close, fires spread, and secondary collapses alter feasible
paths. Decision support systems must therefore combine \emph{geometric} reasoning
(paths and travel time), \emph{combinatorial} reasoning (who is picked up when, under
capacity), \emph{predictive} reasoning (risk and survival proxies), and
\emph{uncertainty-aware} reasoning (ambiguous threat levels interpreted gradually
rather than as brittle thresholds). Piecemeal solutions---e.g., only running Dijkstra
or only training a neural network---do not satisfy a CCP that explicitly demands
\textbf{integration} with traceable behaviour and comparative evaluation.

\subsection{Contributions}
\Aidra{} contributes: (i) a \textbf{reference architecture} tying search, CSP, ML,
and fuzzy modules to one simulation state machine; (ii) a \textbf{working web
deployment} suitable for examiner demonstration; (iii) \textbf{transparent metrics}
(search expansions, CSP backtracks, ML accuracy and macro-F1, KPI summaries); and
(iv) \textbf{documented methodology} bridging theory and the TypeScript codebase,
supporting defence during oral evaluation.

\subsection{CLO mapping (AIC~201)}
This work maps to CLO3 (``appraise real world problems for machine learning based
solutions'') in two ways. First, we justify why \emph{classification-style} risk
labels complement triage CSP objectives and discuss limits of purely synthetic labels.
Second, we show ML not as an endpoint but as a \textit{subscriber} into planning:
model choice alters explanatory priority lines and CSP tie-break cues. Other CLOs
are supported implicitly (problem formulation, team engineering, reproducible
artifacts) through the structured report, repository, and video.

\section{CCP Problem Summary and Formal Sketch}

\subsection{Environment}
Let the map be a grid graph $G=(V,E)$ with $|V|=18\times 18$. Each cell $v$ has type
$T(v)$ (road, safe, fire, collapse, blocked, medical center variants, rescue base),
passability $\phi(v)\in\{0,1\}$, and risk $r(v)\in[0,1]$. Some cells initially host
\textbf{five} victims $\mathcal{V}=\{v_1,\ldots,v_5\}$ with severity $s_i\in$
\{critical,\,moderate,\,minor\}, survival fraction $\sigma_i\in[0,100]$, status
(waiting, en-route, rescued, lost), and estimated time of arrival (ETA) metadata when
moving. Ambulances $A_1,A_2$ occupy cells, maintain ordered assigned victim IDs with
\textbf{capacity two}, and accumulate routes computed by search. Rescue team $R$
may target at most one victim at a time. Ten medical kits appear as aggregate
inventory suitable for narration in reports.

\subsection{State and transitions}
Discrete simulation steps advance global time when autoplay runs; between ticks the
\textbf{paused} engineer may edit coordinates or inject events (after-shock hazard
update, deterministic road block, fire spread seed). Abstract state $S_t$ aggregates
grid fields, actor poses, queues, KPI counters, CSP solution snapshot, fuzzy snapshot
multipliers, ML estimates per victim, and an append-only \textbf{decision log}.
\textbf{Actions} correspond to simulator controls: START triage workflow, TICK,
REPLAN, route refreshes post-assignment, and export of JSON run snapshots.

\subsection{Objectives and metrics}
The user selects among \textit{primary policy emphasis}: minimize travel-related
risk, minimize time-like cost, or a balanced blend---implemented as selectable
objectives that parameterize heuristic and step-cost composition in A* family
methods. Operational metrics surfaced in-client include rescued victim counts,
elapsed mission time (seconds), KPI cards (labels vary by instrumentation), CSP
\textbf{backtrack} counts, frontier and expansion statistics for comparative search,
and classifier \textbf{accuracy} plus \textbf{macro-averaged F1} across three risk
classes. Let $|\mathcal{V}_{\mathrm{resc}}|$ be rescued victims by horizon end:
\begin{equation}
\label{eq:saved}
\mathrm{RescueFraction}=
\frac{|\mathcal{V}_{\mathrm{resc}}|}{|\mathcal{V}|}.
\end{equation}
Average ETA (where defined) summarizes temporal performance; cumulative exposure
along executed routes can be reasoned verbally from logged risk-heavy edges even if
not always shown as one scalar label.

\section{System Architecture}

Figures~\ref{fig:archTwin} and~\ref{fig:archAi} summarize the system at a conceptual
level (operator-facing shell and digital twin, then intelligent modules inside the
coordinator). Export your Mermaid renders to \texttt{figs/architecture-twin.pdf} and
\texttt{figs/architecture-ai.pdf} (vector PDF preferred, width matched to one column).
\textbf{Presentation layer} React components (\texttt{App.tsx}, \texttt{LiveSim.*},
\texttt{SearchTrace.tsx}, \texttt{CspSolver.tsx}, \texttt{MlStudio.tsx},
\texttt{Analytics.tsx}) render shared state---no duplicate ``shadow'' models outside
the engine. \textbf{Simulation kernel} (\texttt{simulationEngine.ts}) validates
events, clones safe state, merges CSP outputs into victim assignments and routes,
applies fuzzy scaling when enabled, invokes ML inference using the retained
evaluation snapshot, and logs human-readable rationales (\texttt{fuzzyLogic.ts},
\texttt{mlRiskPipeline.ts}). \textbf{Search services} encapsulate uninformed /
informed strategies over the same neighbour graph (4-connected unless blocked).
\textbf{CSP solver} emits ambulance victim lists plus team designation and records
\textbf{performance rows} consumed in analytics plots. Tables~\ref{tab:modules}
and~\ref{tab:ui_tabs} summarize code and UX mapping.

\begin{figure}[!t]
  \centering
  \IfFileExists{figs/architecture-twin.pdf}{%
    \includegraphics[width=\columnwidth]{figs/architecture-twin.pdf}%
  }{%
    \fbox{%
      \begin{minipage}[c][0.26\textheight]{0.92\columnwidth}
        \centering
        \vfill
        \footnotesize
        \textbf{Placeholder (a).} Place vector PDF \texttt{figs/architecture-twin.pdf} here
        (Mermaid Fig.~1: operator $\rightarrow$ views $\rightarrow$ control $\rightarrow$
        digital twin $\rightarrow$ outputs).
        \vfill
      \end{minipage}%
    }%
  }
  \caption{Conceptual architecture (part~a): operator, mission views, session control,
  digital twin core (world, agents, victims, plans), and observable outputs.}
  \label{fig:archTwin}
\end{figure}

\begin{figure}[!t]
  \centering
  \IfFileExists{figs/architecture-ai.pdf}{%
    \includegraphics[width=\columnwidth]{figs/architecture-ai.pdf}%
  }{%
    \fbox{%
      \begin{minipage}[c][0.26\textheight]{0.92\columnwidth}
        \centering
        \vfill
        \footnotesize
        \textbf{Placeholder (b).} Place vector PDF \texttt{figs/architecture-ai.pdf} here
        (Mermaid Fig.~2: coordinator $\rightarrow$ search, CSP, ML, fuzzy, local refinement
        $\rightarrow$ updated mission state).
        \vfill
      \end{minipage}%
    }%
  }
  \caption{Conceptual architecture (part~b): intelligent modules invoked by the central
  coordinator and merged back into mission state consumed by part~(a).}
  \label{fig:archAi}
\end{figure}

\begin{table}[!t]
  \caption{Repository modules mapped to grading criteria.}
  \label{tab:modules}
  \centering
  \footnotesize
  \begin{tabular}{@{}p{0.34\linewidth}p{0.62\linewidth}@{}}
    \toprule
    \textbf{Criterion theme} & \textbf{Implementation files / notes} \\
    \midrule
    Search \& heuristics & \texttt{search.ts}, \texttt{heuristics.ts}: BFS/DFS/Greedy/A*, risk/time blend \\
    CSP resource allocation & \texttt{csp.ts}: backtracking, heuristics, counting \\
    Local search refinement & \texttt{localSearch.ts}: hill climbing, simulated annealing \\
    ML appraisal (CLO3) & \texttt{mlRiskPipeline.ts}: synthetic generator, NB/kNN/MLP evaluation \\
    Uncertainty \& fuzzy reasoning & \texttt{fuzzyLogic.ts}: aggregates + rule strengths \\
    World dynamics & \texttt{gridGenerator.ts}: initial disasters; engine mutates hazards \\
    Orchestration \& logging & \texttt{simulationEngine.ts}, UI decision log \\
    \bottomrule
  \end{tabular}
\end{table}

\begin{table}[!t]
  \caption{Major interactive surfaces for evaluation.}
  \label{tab:ui_tabs}
  \centering
  \footnotesize
  \begin{tabular}{@{}p{0.28\linewidth}p{0.68\linewidth}@{}}
    \toprule
    \textbf{Surface} & \textbf{Evaluates / demonstrates} \\
    \midrule
    Live Simulation & Dynamic grid, KPIs, replan, CSP+search coupling, fuzzy toggle \\
    Search Trace & Multi-algorithm comparison, expansions, time, heuristic effects \\
    CSP Solver & Assignment structure + backtracks + explanatory lines \\
    ML Studio & Cross-validated scores, confusion insight, chosen model coupling \\
    Analytics & Consolidated KPI / scenario aggregates for narration \\
    \bottomrule
  \end{tabular}
\end{table}

\section{Detailed Methodology}

\subsection{Neighbour exploration and movement costs}
For each occupiable cell we generate up to four neighbours subject to borders and
passability. Planner \textbf{step cost} merges geographic progression with localized
risk. When fuzzy routing is active, multiplying factors from the fuzzy snapshot scale
risk-sensitive components so identical spatial moves become more expensive near
interpreted hazards or under high organisational urgency---mirroring qualitative
\textbf{staff caution} instructions without rewriting low-level graphs.

\subsection{Search algorithms}
Table~\ref{tab:searchcmp} contrasts theoretical guarantees on uniform-cost grids.

\subsubsection{BFS / DFS}
BFS discovers least-hop routes when edge costs unify; DFS explores deeply with low
memory but may return long detours unsuitable under time KPIs unless depth bounded.

\subsubsection{Greedy best-first}
Expands minimally $f=g+h$ biased toward heuristic alone (implementation-specific tie
break ordering). Useful for responsiveness demonstrations but not globally optimal.

\subsubsection{A* family}
$f(n)=g(n)+h(n)$ with heuristic informed by Euclidean / Manhattan scaffolding plus
risk-weighted lookahead consistent with objectives. Changing objective priority
\textbf{does not duplicate graph code}; it rewires heuristic contributions and expands
comparison plots in Search Trace enabling rubric-aligned quantitative narrative.
In the deployed build the admissible heuristic combines \textbf{Manhattan distance} with
an explicit \textbf{risk penalty} on traversed cells so $h$ remains goal-directed while
penalising exposure.

\subsubsection{Empirical planner comparison (Search Trace)}
On a representative logged run, \textbf{BFS} expanded $167$ nodes, yielded path cost
$55.61$, and cumulative route risk $3.00$. \textbf{DFS} expanded $319$ nodes but produced
a highly suboptimal path (cost $954.02$, risk $53.51$), illustrating depth-first wander
without cost optimality. \textbf{Greedy} minimised search effort ($18$ expansions) with
cost $55.61$ and risk $3.00$---matching BFS on this instance---yet remains a best-first
approximation without global optimality guarantees. \textbf{A*} expanded $76$ nodes,
achieved the lowest path cost ($53.93$) and lowest cumulative risk ($2.71$) while
preserving optimality under the stated heuristic, i.e.\ the best empirical trade-off
between route quality and risk-aware cost accumulation among the four planners.

\subsection{Constraint satisfaction formulation}
\textbf{Decision variables}: assignment of victims to Ambulance~1 roster, Ambulance~2
roster (each length $\leq 2$), and designation of rescue team singleton (or absence).
\textbf{Domains}: feasible victim ID subsets respecting distinctness.
\textbf{Hard constraints}: capacity $\leq 2$ per ambulance, each victim allocated at
most once among ambulances and team feasibility, synergy with medically grounded
\textbf{severity ordering} biases (implementation applies ordering heuristics and
checks). \textbf{Soft integration}: ML-derived risk summaries and fuzzy CSP bump may
shuffle tie resolution among equally consistent partial assignments. The solver uses
\textbf{recursive backtracking} with variable/value ordering enhancements and forwards
illegal combinations; instrumentation exposes \textit{Backtracks} counters for empirical
comparison when toggling heuristic guidance or objective emphasis.
Aggregated for reporting, the UI-facing formulation comprises \textbf{four} primary
assignment variables, \textbf{six} explicit constraints, and combinatorial domains on the
order of \textbf{five} admissible atomic outcomes per orchestrated slot; logged sessions
showed \textbf{all constraints satisfied} with \textbf{zero} violations and
\textbf{zero} backtracks under \textbf{MRV}, \textbf{degree} tie-breaking, and
\textbf{forward checking}, while the produced ordering successfully prioritised
\textbf{critical} victims ahead of moderate/minor cases.

\subsection{Local search overlay}
Beyond tree search for paths, structured assignment or ordering candidates can be
\textbf{energy}-minimized by hill climbing (\textit{deterministic ascend}) or simulated
annealing (\textit{stochastic escapes}) as implemented---useful pedagogically to tie
\textbf{combinatorial optimisation} coursework to CSP outputs.

\subsection{Machine learning pipeline}

\subsubsection{Rationale under data scarcity}
Field EMS labels are scarce; supervised learning nonetheless trains critical thinking if
\textbf{limitations} remain explicit (Section~\ref{sec:discussion}). Synthetic
generators with stratified splitting emulate variance while preserving class balance
manipulation knobs.

\subsubsection{Features and labeling}
Eight continuous features summarise spatial location, surrogate distance-from-base,
aggregated hazard core, fire hint, collapse hint, traffic scalar, weather scalar---
see Table~\ref{tab:feat}. Targets are ordinal-style three-class risk strata.

\begin{table}[!t]
  \caption{Risk-model feature channels (conceptual semantics).}
  \label{tab:feat}
  \centering
  \footnotesize
  \begin{tabular}{@{}clp{0.55\linewidth}@{}}
    \toprule
    \# & Symbol / name & Interpretation \\
    \midrule
    1--2 & $u_r,u_c$ & Normalised row/column coordinates (context) \\
    3 & $d_b$ & Distance-to-base surrogate (supply chain mental model) \\
    4 & $h_{\mathrm{core}}$ & Hazard core aggregate before incident modifiers \\
    5--6 & $f^+, c^+$ & Fire tendency / structural collapse tendency hints \\
    7--8 & $\tau,w$ & Traffic \& weather perturbation scalars \\
    \bottomrule
  \end{tabular}
\end{table}

\subsubsection{Models and optimisation}
\textbf{k-NN}: Euclidean distance across z-scored features with cross-validated
$k\in\{1,2,3,5,7,10\}$; the deployment snapshot retained $k=5$ as the best trade-off.
\textbf{Gaussian NB}: class-conditional Gaussian with variance floor $\varepsilon$.
\textbf{MLP}: softmax classifier with architecture $8\to 16\to 3$, ReLU activations, and
SGD optimisation (learning rate $0.08$, $80$ epochs). It attained the strongest hold-out
scores (Table~\ref{tab:metrics}), consistent with capturing \textbf{nonlinear}
interactions among hazard proximity, severity, and blockage likelihood. Post-hoc
importance diagnostics from the evaluation UI rank \texttt{distance\_to\_hazard} at
$0.82$, \texttt{victim\_severity} at $0.78$, and \texttt{road\_blockage\_prob} at $0.71$,
reinforcing face-valid emphasis on spatial threat, triage acuity, and infrastructure
fragility.

\begin{table}[!ht]
  \caption{ML dataset and training hyperparameters (values from deployed evaluation run).}
  \label{tab:ml_hyp}
  \centering
  \footnotesize
  \begin{tabular}{@{}ll@{}}
    \toprule
    \textbf{Item} & \textbf{Setting} \\
    \midrule
    Samples & $N=500$ synthetic \\
    Train/test split & 80/20 stratified \\
    Feature dimension & 8 \\
    Classes & 3 risk classes (Low / Medium / High) \\
    Best k-NN & $k=5$ \\
    MLP architecture & $8\to 16\to 3$, ReLU \\
    MLP optimisation & SGD, LR $=0.08$, $80$ epochs \\
    Reported scores & Accuracy, Macro-F1 (Table~\ref{tab:metrics}) \\
    \bottomrule
  \end{tabular}
\end{table}

\subsection{Fuzzy uncertainty layer}
Three \textbf{crispified} aggregates feed triangular memberships: hazard $h$ from fire,
collapse, blocked density, smooth road-risk mean; urgency $u$ derived from inverted
waiting survival fractions; uncertainty $x$ elevating blockage + residual variance.
\textbf{Rules} aggregate via minimum t-norm strengths and drive consequent singleton
weights adjusting (i)~risk step multiplier, (ii)~heuristic risk weight multiplier,
(iii)~CSP micro-priority elevation; first fired rule exposition appears in snapshots
feeding operator explanations in the fuzzy panel alongside crisp numeric echoes.

\section{Operational Loop and Experimental Protocol}

\subsection{Triage, simulation, replanning}
\textbf{START} establishes consistent initial assignments when needed. \textbf{Autoplay}
increments time with selectable speed tiers; paused operation supports inspect-and-edit.
\textbf{Replan} re-solves CSP under current degradation and rebinds sequential routes---
essential rubric storyline: ``world changed $\rightarrow$ plan invalid $\rightarrow$
transparent revision with logging.''

\subsection{Suggested repeatable trials}
Controlled narrative for table population: (1)~baseline static grid, greedy vs A*,
record expansions \& path cost; (2)~same baseline, fuzzy off vs on, annotate multiplier
movement; (3)~stress test with mid-run block injecting detour latency; (4)~ML sweep
changing active model observing stability of CSP tie behaviour; each trial export JSON via
\textbf{Export} button for appendix material if graders allow uploads beyond six pages online.

\subsection{Where to obtain values for empirical tables (step-by-step)}
Table~\ref{tab:fill_guide} is your \textbf{cheat sheet}: each row in Table~\ref{tab:metrics}
maps to a concrete screen in the running \Aidra{} app (navbar tabs + side panels). Open
\texttt{npm run dev}, use a fixed window size when screenshotting, and copy numbers after
each action so the report matches what the examiner can reproduce.

\begin{table*}[!t]
  \caption{Guide: where to read each number before typing it into Table~\ref{tab:metrics}.}
  \label{tab:fill_guide}
  \centering
  \footnotesize
  \begin{tabular}{@{}p{0.26\linewidth}p{0.70\linewidth}@{}}
    \toprule
    \textbf{Report row / metric} & \textbf{In the application (exact path)} \\
    \midrule
    ML k-NN / NB / MLP Accuracy \& Macro-F1 &
    Navbar $\rightarrow$ \textbf{ML Studio}. If metrics show ``---'', click
    \textbf{Run Training \& Evaluation} so \texttt{mlEvalSnapshot} populates.
    Each model card shows \textit{Acc} (\%) and \textit{F1} (two decimals). You can
    cross-check the same three lines on \textbf{Live Sim}: right column \textbf{Fuzzy/ML}
    strip (narrow cards for kNN, NB, MLP). \\
    Greedy vs A* nodes expanded, time, path cost &
    Navbar $\rightarrow$ \textbf{Search Trace}. Run search for the configured algorithm
    (UI triggers \texttt{runSearchForAlgorithm}). The \textbf{comparison table} lists
    \textit{Nodes expanded}, \textit{Time (ms)}, \textit{Path cost}, and optimality flags
    per algorithm row. Pick Greedy and A* rows for your narrative table. \\
    CSP backtracks (scenario A / B) &
    Navbar $\rightarrow$ \textbf{CSP Solver}. Click \textbf{Run Solver}. The decision
    \textbf{Log} on Live Sim (or CSP narrative text) includes \texttt{backtracks:\,N}
    after a successful solve. Run twice: (A) fresh
    simulation after reset; (B) after placing a \textbf{Block} or \textbf{Fire} then
    \textbf{Replan} so the assignment problem tightens---record both $N$ values. \\
    Rescued vs total victims &
    \textbf{Live Sim}: watch victim \textbf{Status} in the right \textbf{Victims} table
    (\texttt{rescued} vs others), or \textbf{Analytics} tab for aggregated scenario row
    text that includes saved count once \texttt{allAlgoComparisons} is non-empty. KPI
    cards at the bottom of the live layout also list instructor-facing labels. \\
    Avg ETA (optional) &
    \textbf{Live Sim} $\rightarrow$ right \textbf{Victims} table, \textbf{ETA} column
    (minutes). Average manually for moderate rows or cite a single representative run. \\
    Fuzzy multiplier range &
    \textbf{Live Sim}: enable fuzzy in \textbf{AI Configuration} (left panel dropdown).
    Read \textbf{s$\times$}, \textbf{h$\times$}, \textbf{csp+} lines in the fuzzy card
    (same strip as ML). Toggle fuzzy OFF, note neutral values, toggle ON after a shock
    and record min/max you observed across two replans. \\
    Path optimality, risk exposure, replan count, avg rescue time &
    Navbar $\rightarrow$ \textbf{Analytics}. After running searches from \textbf{Search Trace},
    the \textbf{live scenario} row shows \textit{victims saved} as \texttt{X/5},
    \textit{avg time}, rounded \textit{risk exposure}, \textit{optimality} (or ``---''
    until costs exist), and \textit{replan} count---these map directly to
    \texttt{liveScenarioRow} in code. \\
    \bottomrule
  \end{tabular}
\end{table*}

\section{Related Work and Positioning}
Classical planning literature separates \emph{pathfinding} (shortest paths, A*~\cite{hart1968})
from \emph{assignment} and scheduling (often modeled as CSP or mixed integer programs).
Recent disaster informatics layers machine learning on top of GIS feeds for damage
assessment; our scope is deliberately smaller---a self-contained grid---but preserves
the \textbf{integration challenge}: learned scores must not silently override hard
capacity constraints; fuzzy modifiers must remain bounded so optimality arguments for
A* variants remain discussable in viva. Russell and Norvig~\cite{russell2021} supply the
common vocabulary (search, CSP, uncertainty) we reference when explaining design trade-offs.

\section{Computational Complexity and Scalability Notes}
Let $n=|V|$ be open cells and $d$ maximum branching (four on a grid). Uninformed BFS/DFS
worst-case time is $O(n)$ expansions for reachability-style goals on implicit graphs with
duplicate suppression. A* with a consistent heuristic explores $O(n)$ in the worst case
but typically far fewer when $h$ is informative. The CSP assignment layer has worst-case
exponential backtracking in the number of victims; with only five victims and small
domains the observed backtrack counts remain human-interpretable---use them as empirical
complexity proxies rather than asymptotic proofs. ML training is $O(N^2)$ per epoch for
naive k-NN scoring on $N{=}500$ samples (acceptable in-browser) and linear-ish per epoch
for the compact MLP.

\begin{table}[!ht]
  \caption{Asymptotic and practical complexity sketch (qualitative).}
  \label{tab:complexity}
  \centering
  \footnotesize
  \begin{tabular}{@{}p{0.34\linewidth}p{0.60\linewidth}@{}}
    \toprule
    \textbf{Module} & \textbf{Typical cost driver in \Aidra} \\
    \midrule
    BFS / DFS & Frontier volume; duplicate detection hash \\
    Greedy / A* & Priority-queue ops; heuristic evaluation per edge \\
    CSP & Backtracking depth; constraint checks per partial assignment \\
    k-NN inference & Distance to all training points per victim feature vector \\
    MLP inference & Two matrix-vector products (tiny) \\
    Fuzzification & O(cells) scan for aggregates + fixed rule count \\
    \bottomrule
  \end{tabular}
\end{table}

\section{Formal A* Cost Model (summary)}
Let $g(n)$ accumulate nonnegative step costs along edges from the start cell to node $n$,
and let $h(n)$ estimate remaining cost to a goal cell (medical center or intermediate target).
A* orders OPEN by $f(n)=g(n)+h(n)$. Risk-aware settings add weighted cell-risk terms inside
$g$ and optionally inside $h$; fuzzy ON multiplies the risk-sensitive portion so identical
topologies yield different $f$-orderings---this is the quantitative hook for comparing
\textbf{Crisp} vs \textbf{Fuzzy} runs in the Search Trace table without altering underlying
connectivity. Empirically, enabling fuzzy reasoning alongside A* lowered aggregate risk
exposure relative to uninformed baselines (e.g., DFS) on the logged benchmark, aligning
operator-facing multiplier telemetry (Table~\ref{tab:metrics}) with safer route envelopes.

\section{Software Engineering and Verification}
The project uses Vite, React~18, and TypeScript for static typing across engine and UI.
Vitest unit tests (\texttt{src/utils/*.test.ts}) guard analytics helpers; interactive
regression remains manual via the checklist in Table~\ref{tab:qa}. Continuous integration
(if enabled in your fork) should run \texttt{npm test} and \texttt{npm run build} before
tagging submission releases.

\begin{table}[!ht]
  \caption{Manual QA checklist before recording the demo or printing numbers into tables.}
  \label{tab:qa}
  \centering
  \footnotesize
  \begin{tabular}{@{}cp{0.82\linewidth}@{}}
    \toprule
    $\checkmark$ & Reset simulation; confirm five victims visible on grid. \\
    $\checkmark$ & Run ML evaluation once; confirm no ``---'' in ML Studio metrics. \\
    $\checkmark$ & Run all search algorithms from Search Trace; confirm comparison rows fill. \\
    $\checkmark$ & CSP Solver $\rightarrow$ \textbf{Run Solver}; confirm backtrack count in log text. \\
    $\checkmark$ & Toggle fuzzy; replan; confirm multiplier lines change. \\
    $\checkmark$ & Export JSON snapshot; archive file name with date for your appendix zip. \\
    \bottomrule
  \end{tabular}
\end{table}

\section{Self-assessment vs.\ course rubric (Assignment~3)}
Table~\ref{tab:rubric} links examiner criteria to evidence you should point to during
defence (not a claim of marks---a coverage map).

\begin{table}[!ht]
  \caption{Rubric alignment (5 marks total): where evidence lives.}
  \label{tab:rubric}
  \centering
  \footnotesize
  \begin{tabular}{@{}p{0.14\linewidth}p{0.18\linewidth}p{0.60\linewidth}@{}}
    \toprule
    \textbf{Wt.} & \textbf{Criterion} & \textbf{Evidence in \Aidra} \\
    \midrule
    1.5 & Design / architecture / implementation &
    Architecture PDF + file map tables + live modular tabs (Search/CSP/ML/Fuzzy). \\
    2.0 & Results / evaluation / technical depth &
    Filled Table~\ref{tab:metrics}, Search Trace comparisons, ML Studio metrics, CSP backtracks, discussion paragraphs. \\
    1.5 & GitHub + LinkedIn video &
    Public repo README with build steps + embedded or linked 2-minute video URL. \\
    \bottomrule
  \end{tabular}
\end{table}

\section{Evaluation Metrics Snapshot}
Table~\ref{tab:metrics} is the primary quantitative dashboard fill-in template.
Table~\ref{tab:kpi} enumerates KPI conceptual definitions aligned with narration.

\subsection{Formal definitions: macro-F1 and path optimality}
For model $m$, let $P_c,R_c,F_c$ be precision, recall, and F1 for class
$c\in\{L,M,H\}$. Macro-F1 averages class-wise F1 with equal weight:
\begin{equation}
\mathrm{F1}_{\mathrm{macro}}^{(m)}=\frac{1}{3}\sum_{c\in\{L,M,H\}} F_c^{(m)}.
\end{equation}
The ML Studio master table and BottomBar strip surface the aggregated macro-F1 from
the evaluation snapshot. For path optimality, given nonnegative path costs
$\mathrm{cost}(a)$ for each algorithm $a$ that produced a route in the same scenario,
let $\mathrm{cost}^\star=\min_a \mathrm{cost}(a)$. The selected algorithm $s$ has score
$\min\bigl(1,\mathrm{cost}^\star/\mathrm{cost}(s)\bigr)$ in $[0,1]$
(\texttt{pathOptimalityScore} in \texttt{analyticsLive.ts}). Unity means the chosen
planner matched the cheapest observed peer on that run.

\begin{table}[!t]
  \caption{Empirical readings from the deployed \Aidra{} application (representative session).}
  \label{tab:metrics}
  \centering
  \footnotesize
  \begin{tabular}{@{}p{0.42\linewidth}cc@{}}
    \toprule
    \textbf{Observable} & \textbf{Value(s)} & \textbf{Capture location} \\
    \midrule
    ML k-NN Accuracy / Macro-F1 & 81.2\% / 0.745 & ML Studio \\
    ML Naive Bayes Acc / Macro-F1 & 80.2\% / 0.714 & ML Studio \\
    MLP Acc / Macro-F1 & 98.0\% / 0.973 & ML Studio \\
    Greedy expansions vs A* expansions (same pair) & 18 / 76 & Search Trace \\
    CSP backtracks (scenario A/B) & 0 / 0 & CSP solver panel \\
    Rescued vs total victims & 5 / 5 & Live KPI HUD \\
    Avg ETA (moderate tier subset) optional & 8.1 min & Victims panel \\
    Fuzzy multiplier range observed & approx.\ $1.0\times$\,--\,$1.8\times$ (risk weight) & Live fuzzy card \\
    \bottomrule
  \end{tabular}
\end{table}

\begin{table}[!ht]
  \caption{Narrative KPI definitions (adapt wording to examiner rubric wording).}
  \label{tab:kpi}
  \centering
  \footnotesize
  \begin{tabular}{@{}p{0.32\linewidth}p{0.64\linewidth}@{}}
    \toprule
    \textbf{KPI} & \textbf{Meaning \& interpretation} \\
    \midrule
    Rescue throughput & Ratio (\ref{eq:saved}); failure analysis if $<1$ cites survival decay \\
    Search effort & Expanded nodes \& wall clock proxy efficiency trade-off \\
    CSP hardness & Larger backtracks $\Rightarrow$ tighter assignment coupling \\
    Classifier calibration & Accuracy vs macro-F1 gap highlights class imbalance resilience \\
    Fuzzy coherence & Multipliers should correlate textual rule preview with KPI shifts \\
    \bottomrule
  \end{tabular}
\end{table}

\subsection{Representative analytics scenarios (S1--S4)}
Table~\ref{tab:metrics} summarises global ML and planner telemetry; the Analytics tab
further supports \emph{scenario-style} contrasts when operators pair algorithms with
model selectors. \textbf{S1} (BFS + kNN): victims saved $3/5$, risk score $78$,
optimality $0.71$. \textbf{S2} (DFS + NB): saved $3/5$, risk $82$, optimality $0.64$.
\textbf{S3} (Greedy + kNN): saved $4/5$, risk $54$, optimality $0.79$. \textbf{S4}
(A* + fuzzy + MLP): saved $5/5$, risk $38$, optimality $0.94$. Taken together, the
strongest aggregate outcome arises when informed search, fuzzy cost shaping, and the
best-calibrated risk model co-apply; in our logged campaign,
\textit{A* combined with fuzzy reasoning and the MLP risk model achieved the best overall balance between rescue efficiency and risk reduction.}

\begin{table}[!ht]
  \caption{Comparative search algorithm summary (theory-aligned).}
  \label{tab:searchcmp}
  \centering
  \footnotesize
  \begin{tabular}{@{}lccp{0.40\linewidth}@{}}
    \toprule
    \textbf{Alg.} & \textbf{Optim.} & \textbf{Complete} & \textbf{Remarks for demo script} \\
    \midrule
    BFS & Yes (uniform) & Yes & Show optimality baseline on simple pair \\
    DFS & No & Yes (finite) & Contrast zig-zag artefacts \\
    Greedy & No & Yes & Fast first answer; quantify suboptimality \\
    A* & If $h$ admiss.\ & Yes & Stress objective toggling narrative \\
    \bottomrule
  \end{tabular}
\end{table}

\subsection{CSP backtracking skeleton (informative pseudocode)}
\begin{algorithm}[!t]
\caption{Informative CSP backtracking overview}
\label{alg:csp}
\begin{algorithmic}[1]
\Procedure{CSPAssign}{partial, domains}
    \If{\textit{complete}(partial)}
        \State \Return partial
    \EndIf
    \State $X \gets \textsc{SelectVar}(\mathit{domains})$
      \Comment{variable ordering heuristic}
    \For{each value $v$ in \textsc{OrderedDomain}$(X)$}
        \If{\textit{consistent}$(X{=}v$, partial)}
            \State $\mathit{domains}' \gets$ propagate / prune successors
            \If{no empty domain}
                \State $res \gets$ \Call{CSPAssign}{partial $\cup\{X{=}v\}$, $\mathit{domains}'$}
                \If{$res$ is feasible}
                    \State \Return $res$
                \Else
                    \State increment backtrack counter
                \EndIf
            \EndIf
        \EndIf
    \EndFor
    \State \Return failure
\EndProcedure
\end{algorithmic}
\end{algorithm}

\begin{table*}[!t]
  \caption{Optional extended lab book: duplicate block per scripted experiment (print blank PDF row if unused).}
  \label{tab:labbook}
  \centering
  \scriptsize
  \begin{tabular}{@{}p{0.06\linewidth}p{0.18\linewidth}p{0.12\linewidth}p{0.12\linewidth}p{0.12\linewidth}p{0.32\linewidth}@{}}
    \toprule
    ID & Scripted actions (short) & Saved & Risk & Opt. & Notes (fuzzy on/off, algorithm) \\
    \midrule
    R1 & S1: BFS + kNN & 3/5 & 78 & 0.71 & Analytics snapshot \\
    R2 & S2: DFS + NB & 3/5 & 82 & 0.64 & Analytics snapshot \\
    R3 & S3: Greedy + kNN & 4/5 & 54 & 0.79 & Analytics snapshot \\
    R4 & S4: A* + fuzzy + MLP & 5/5 & 38 & 0.94 & Analytics snapshot \\
    R5 & Change ML model $\rightarrow$ Replan & & & & \\
    \bottomrule
  \end{tabular}
\end{table*}

\section{Discussion}\label{sec:discussion}

\subsection{Interpretation playbook}
The MLP dominated the synthetic risk triage task ($98.0\%$ accuracy, macro-F1 $0.973$;
Table~\ref{tab:metrics}), outperforming k-NN and Naive Bayes and supporting its use as the
default high-fidelity estimator in scenario~S4. Greedy search minimised expansions
($18$) but, as theory predicts, forfeits optimality certificates; A* required more
work ($76$ expansions) yet delivered the lowest path cost ($53.93$) and lowest logged
route risk ($2.71$), outperforming BFS on cost while avoiding DFS's extreme detour
(cost $954.02$, risk $53.51$). Fuzzy-weighted A* runs further depressed aggregate exposure
versus uninformed routing, matching the observed heuristic risk multiplier band
(Table~\ref{tab:metrics}). CSP telemetry remained benign in both stress snapshots
($0/0$ backtracks) with MRV, degree, and forward checking, indicating that the current
five-victim abstraction stays within the solver's easy regime while still demonstrating
full constraint satisfaction and critical-first ordering.

\subsection{Limitations}
Synthetic labels omit spatial autocorrelation; fuzzification flattens fine-grained
weather/wind cues. Interactive runs require scripted seeds for bit-identical replay.

\subsection{Ethics \& responsible AI framing}
Demonstrations emphasize decision support (\textit{assistive}), not autonomous vehicle
authority. Exported logs should anonymise if ever replaying real municipality trials.

\subsection{Threats to validity (study design)}
\textbf{Construct validity:} synthetic ML labels may not correlate with real triage
outcomes---treat classifiers as \emph{pedagogical} rankers. \textbf{Internal validity:}
interactive operator order affects logs; document a scripted click path. \textbf{External
validity:} grid topology omits multi-floor hospitals and traffic dynamics. Acknowledging
these threats strengthens viva defence rather than weakening it.

\subsection{Group work split (template)}
Replace with your actual division: Member~A owns search/heuristic integration and Search
Trace validation scripts; Member~B owns CSP instrumentation and report tables; joint
sessions cover fuzzy tuning and video storyboard. Version-control merge commits should
reflect this split superficially for academic integrity interviews.

\section{Glossary of recurring symbols}
\begin{table}[!ht]
  \caption{Notation quick reference (non-exhaustive).}
  \label{tab:glossary}
  \centering
  \footnotesize
  \begin{tabular}{@{}cp{0.82\linewidth}@{}}
    \toprule
    Symbol & Meaning in this report / codebase \\
    \midrule
    $G=(V,E)$ & Grid graph; $V$ cells, $E$ cardinal adjacency when passable \\
    $\phi(v)$ & Passability indicator for cell $v$ \\
    $r(v)$ & Scalar risk in $[0,1]$ stored on road-like cells \\
    $\sigma_i$ & Survival percentage of victim $i$ \\
    $g,h,f$ & A* accumulated cost, heuristic, combined priority \\
    $h,u,x$ & Fuzzy crisp inputs: hazard, urgency, uncertainty aggregates \\
    \bottomrule
  \end{tabular}
\end{table}

\section{Demonstration \& Submission Checklist}

\subsection{Two-minute LinkedIn/GitHub narrative arc}
\textbf{(0--20\,s)} Problem statement \& objectives; \textbf{(20--60\,s)} Live map CSP
solve + single replan event; \textbf{(60--100\,s)} Search Trace contrasting two
planners with numeric callouts; \textbf{(100--120\,s)} ML Studio classifier evidence +
fuzzy ON/OFF micro comparison + closing roadmap.

\section*{Artifacts (URLs)}
\textit{(Replace placeholders prior to compilation.)}
\begin{itemize}
\item GitHub repo: \url{https://github.com/YOUR_ORG/YOUR_REPO}
\item Demo / LinkedIn: \url{https://www.linkedin.com/posts/YOUR_POST}
\item Optional JSON exports: attach under \texttt{artifacts/}\ in repository.
\end{itemize}

\section{Conclusion}
\Aidra{} operationalises the CCP as an integrated socio-technical AI stack where
\textbf{representation} (grid \& actors), \textbf{reasoning kernels} (search + CSP +
local search), \textbf{inductive modules} (ML), and \textbf{graduated uncertainty}
(fuzzy) cohabit within one audited state trail. The logged deployment corroborates the
design intent: MLP-guided estimates, zero-backtrack CSP assignments, and A*+fuzzy routing
jointly maximise rescue completion ($5/5$) while minimising measured risk and retaining
optimality guarantees absent in greedy or DFS baselines. Expanded tables and methodological
discussion above support a six-page IEEE PDF once vector architecture art and curated
screenshots are merged; trim optional prose if fractional overflow persists.

\vspace{0.5em}
\noindent\textbf{Acknowledgment.} Course instructor Dr.\ Arshad Farhad.

\balance
\begin{thebibliography}{99}
\bibitem{russell2021}
S.~Russell and P.~Norvig, \emph{Artificial Intelligence: A Modern Approach}, 4th ed.\ Pearson, 2021.
\bibitem{hart1968}
P.~E. Hart, N.~J. Nilsson, and B.~Raphael, ``A formal basis for the heuristic determination of minimum cost paths,'' \emph{IEEE Trans.\ Systems Science and Cybernetics}, vol.~4, no.~2, pp.~100--107, 1968.
\bibitem{zadeh1965}
L.~A.~Zadeh, ``Fuzzy sets,'' \emph{Information and Control}, vol.~8, no.~3, pp.~338--353, 1965.
\end{thebibliography}

\end{document}
